\chapter{Kesimpulan dan Saran}

\section{Kesimpulan}

Berdasarkan hasil penelitian dan pembahasan pada Bab~4 mengenai evaluasi strategi \textit{retrieval} dalam sistem \textit{Retrieval-Augmented Generation} (RAG) untuk pemetaan kontrol \textit{Statement of Applicability} (SoA) ISO/IEC 27001:2022, dapat ditarik kesimpulan sebagai berikut.

\begin{enumerate}
	\item \textit{Hybrid retrieval} dengan proporsi \textit{dense}=0,50:\textit{sparse}=0,50 menghasilkan performa terbaik pada Recall@3=0,537, MAP@3=0,467, dan NDCG@3=0,575. Perbedaan performa antar konfigurasi terbukti signifikan secara statistik ($p < 0{,}001$), yang mengonfirmasi bahwa pemetaan kontrol ISO 27001 membutuhkan kombinasi pencocokan semantik dan leksikal.

	\item \textit{Candidate pool} berukuran 5 menghasilkan performa terbaik pada seluruh metrik (Hit@3=0,857, Recall@3=0,468, MAP@3=0,387, NDCG@3=0,485) dengan latensi terendah (5,91 detik). Penambahan ukuran \textit{pool} justru menurunkan performa secara monoton karena memperbanyak kandidat yang kurang relevan pada tahap \textit{reranking} ($p < 0{,}003$).

	\item \textit{Score fusion} dengan proporsi Hybrid:CE=0,75:0,25 menghasilkan performa terbaik secara keseluruhan (Recall@3=0,541, MAP@3=0,470, NDCG@3=0,579). \textit{Hybrid retrieval} dan \textit{cross-encoder} bersifat komplementer: \textit{hybrid} unggul pada \textit{query} dengan terminologi teknis spesifik (36 dari 77 \textit{query}), sementara \textit{cross-encoder} memberikan nilai tambah pada \textit{query} yang membutuhkan pemahaman kontekstual (15 dari 77 \textit{query}).
\end{enumerate}

Secara keseluruhan, konfigurasi optimal yang teridentifikasi adalah \textit{hybrid retrieval} dengan proporsi \textit{dense}=0,50:\textit{sparse}=0,50, \textit{candidate pool} berukuran 5, dan \textit{score fusion} dengan proporsi Hybrid:CE=0,75:0,25, yang mencapai keseimbangan terbaik antara kualitas \textit{retrieval} dan efisiensi komputasi.

\section{Saran}

Berdasarkan hasil penelitian yang telah dilakukan, beberapa saran untuk penelitian selanjutnya adalah sebagai berikut.

\begin{enumerate}
	\item \textbf{Evaluasi pada domain dan organisasi yang berbeda.} Konfigurasi optimal yang ditemukan perlu diuji pada dokumen kebijakan organisasi lain dengan karakteristik terminologi dan ruang lingkup kontrol yang berbeda untuk memvalidasi generalisasi temuan.

	\item \textbf{Eksplorasi granularitas parameter yang lebih spesifik.} Penelitian ini menguji pembobotan dengan kelipatan 0,25 pada tiga sumbu desain. Penelitian selanjutnya dapat menerapkan \textit{grid search} dengan resolusi lebih halus (kelipatan 0,05 atau 0,10), serta mengevaluasi metode \textit{reranking} alternatif seperti \textit{listwise reranker} atau model berbasis LLM.

	\item \textbf{Analisis granular pada domain kontrol yang kurang terwakili.} Dataset penelitian ini didominasi domain A.5 dan A.8 (96,5\%), sementara A.6 dan A.7 sangat terbatas. Penelitian selanjutnya perlu mengonstruksi dataset dengan cakupan lebih merata pada keempat domain untuk memvalidasi konsistensi konfigurasi optimal.

	\item \textbf{Evaluasi \textit{pipeline} RAG secara menyeluruh.} Penelitian selanjutnya perlu mengevaluasi \textit{pipeline} RAG secara \textit{end-to-end}, meliputi tahap \textit{pre-retrieval} (\textit{document chunking}), \textit{retrieval} dengan konfigurasi optimal, serta tahap \textit{generation} dengan LLM untuk mengevaluasi \textit{faithfulness} dan utilitas rekomendasi kontrol.

	\item \textbf{Validasi hasil oleh praktisi profesional bersertifikat.} Pemetaan manual kontrol Annex pada penelitian ini dilakukan oleh peneliti. Penelitian selanjutnya perlu melibatkan auditor bersertifikat ISO 27001 untuk memvalidasi akurasi \textit{ground truth} dan relevansi hasil pemetaan.

	\item \textbf{Penerapan uji \textit{post-hoc} untuk perbandingan berpasangan.} Penelitian ini menggunakan uji Friedman secara keseluruhan namun belum dilanjutkan uji \textit{post-hoc} seperti Wilcoxon dengan koreksi Bonferroni untuk mengidentifikasi pasangan konfigurasi mana yang benar-benar berbeda secara signifikan.

	\item \textbf{Klasifikasi otomatis jenis \textit{query} berdasarkan karakteristik \textit{retrieval}.} \textit{Hybrid retrieval} dan \textit{cross-encoder} bersifat komplementer bergantung pada karakteristik \textit{query}. Penelitian selanjutnya perlu mengembangkan skema klasifikasi \textit{query} dan mekanisme \textit{adaptive retrieval} yang menyesuaikan pembobotan secara dinamis per \textit{query}.

	\item \textbf{Eksplorasi desain basis pengetahuan dan pemilihan komponen model.} Penelitian selanjutnya perlu menguji variasi desain basis pengetahuan (struktur atribut, cakupan informasi, representasi teks) serta mengeksplorasi kombinasi komponen model (\textit{sparse retrieval}, \textit{dense embedding}, \textit{reranker}) yang lebih efektif pada domain pemetaan kontrol ISO/IEC 27001.
\end{enumerate}
